% Partie : méthodes
% Chapitre : induction mathématique
% Section : récurrence forte

\documentclass[../../../include/open-logic-section]{subfiles}

\begin{document}
\olfileid[fr]{mth}{ind}{str}

\olsection{Récurrence forte}

Dans le principe de récurrence
étudié ci-dessus, nous
démontrons $P(0)$ et
l'implication
«~si $P(k)$, alors
$P(k+1)$~». Pour cette
seconde étape, nous
supposons $P(k)$ et
nous nous en servons
pour démontrer $P(k+1)$.
De manière équivalente,
pour $k\ge1$, nous
pourrions supposer
$P(k-1)$ et en
déduire $P(k)$.
L'essentiel est de
pouvoir passer d'un
nombre quelconque
à son successeur :
démontrer l'énoncé
pour chaque nombre
en supposant qu'il
vaut pour son
prédécesseur.

Il existe une variante
du principe de récurrence
dans laquelle, pour
établir $P(k)$,
nous ne supposons pas
seulement que la
propriété vaut pour
le prédécesseur
$k-1$, mais pour
tous les
nombres inférieurs
à~$k$. Cette méthode
établit également
$P(n)$ pour tout
$n \in \Nat$.
En effet, une fois
$P(0)$ établi,
la propriété vaut
pour tous les
nombres inférieurs
à~$1$. Si nous
avons aussi montré
que $P(k)$ découle
de $P(l)$ pour
tout $l<k$,
appliquons ce
résultat à $k=1$ :
nous obtenons~$P(1)$.
Nous avons alors
$P(0)$ et $P(1)$,
soit $P(l)$ pour
tout $l<2$.
La même implication,
appliquée à $k=2$,
donne~$P(2)$,
et ainsi de suite.

En réalité, si nous
pouvons établir la
condition générale
«~pour tout~$k$,
si $P(l)$ vaut
pour tout $l<k$,
alors $P(k)$~»,
il n'est plus
nécessaire de
démontrer $P(0)$
séparément : elle
en découle déjà.
Rappelons qu'un
énoncé universel
comme «~pour tout
$l<k$, $P(l)$~»
est vrai lorsqu'il
n'existe aucun
$l<k$. C'est
un cas de vérité
par vacuité d'un
énoncé universel :
«~tous les $A$
sont des $B$~»
est vrai s'il
n'existe aucun~$A$ ;
$\lforall[x][(!A(x) \lif !B(x))]$
est vrai si aucun
$x$ ne satisfait~$!A(x)$.
Ici, la version
formalisée est
$\lforall[l][(l < k \lif P(l))]$ ;
elle est vraie
s'il n'existe
aucun $l<k$.
Or c'est précisément
le cas pour $k=0$ :
il n'existe aucun
$l<0$. L'énoncé
«~pour tout
$l<0$, $P(l)$~»
est donc vrai,
quelle que soit~$P$.\footnote{
Correction de l'édition française :
à cet endroit, la source anglaise
écrit «~$P(0)$~» au lieu de
«~$P(l)$~» dans l'hypothèse
universelle ; la vacuité
demeure, mais le schéma
formel porte sur~$P(l)$.}
La démonstration
de «~si $P(l)$
vaut pour tout
$l<k$, alors
$P(k)$~» établit
donc automatiquement
$P(0)$.

Cette variante est
utile lorsqu'on
ne peut pas établir
$P(k)$ à partir
du seul cas $k-1$,
mais qu'on a besoin
de supposer la
propriété pour
un ou plusieurs
nombres $l<k$.

\end{document}
